\section{Degrees of zero cycles}
\label{section-degree-zero-cycles}

\noindent
We start defining the degree of a zero cycle on a proper scheme over a field.
One approach is to define it directly as in
Lemma \ref{lemma-spell-out-degree-zero-cycle} and then show
it is well defined by
Lemma \ref{lemma-curve-principal-divisor}.
Instead we define it as follows.

\begin{definition}
\label{definition-degree-zero-cycle}
Let $k$ be a field (Example \ref{example-field}). Let $p : X \to \Spec(k)$
be proper. The {\it degree of a zero cycle} on $X$ is given by proper
pushforward
$$
p_* : \CH_0(X) \to \CH_0(\Spec(k))
$$
(Lemma \ref{lemma-proper-pushforward-rational-equivalence})
combined with the natural isomorphism $\CH_0(\Spec(k)) = \mathbf{Z}$
which maps $[\Spec(k)]$ to $1$. Notation: $\deg(\alpha)$.
\end{definition}

\noindent
Let us spell this out further.

\begin{lemma}
\label{lemma-spell-out-degree-zero-cycle}
Let $k$ be a field. Let $X$ be proper over $k$. Let $\alpha = \sum n_i[Z_i]$
be in $Z_0(X)$. Then
$$
\deg(\alpha) = \sum n_i\deg(Z_i)
$$
where $\deg(Z_i)$ is the degree of $Z_i \to \Spec(k)$, i.e.,
$\deg(Z_i) = \dim_k \Gamma(Z_i, \mathcal{O}_{Z_i})$.
\end{lemma}

\begin{proof}
This is the definition of proper pushforward
(Definition \ref{definition-proper-pushforward}).
\end{proof}

\noindent
Next, we make the connection with degrees of vector bundles
over $1$-dimensional proper schemes over fields as defined in
Varieties, Section \ref{varieties-section-divisors-curves}.

\begin{lemma}
\label{lemma-degree-vector-bundle}
Let $k$ be a field. Let $X$ be a proper scheme over $k$ of dimension $\leq 1$.
Let $\mathcal{E}$ be a finite locally free $\mathcal{O}_X$-module of constant
rank. Then
$$
\deg(\mathcal{E}) = \deg(c_1(\mathcal{E}) \cap [X]_1)
$$
where the left hand side is defined in
Varieties, Definition \ref{varieties-definition-degree-invertible-sheaf}.
\end{lemma}

\begin{proof}
Let $C_i \subset X$, $i = 1, \ldots, t$ be the irreducible components
of dimension $1$ with reduced induced scheme structure and let $m_i$ be the
multiplicity of $C_i$ in $X$. Then $[X]_1 = \sum m_i[C_i]$ and
$c_1(\mathcal{E}) \cap [X]_1$ is the sum of the pushforwards of the cycles
$m_i c_1(\mathcal{E}|_{C_i}) \cap [C_i]$. Since we have a similar decomposition
of the degree of $\mathcal{E}$ by
Varieties, Lemma \ref{varieties-lemma-degree-in-terms-of-components}
it suffices to prove the lemma in case $X$ is a proper curve over $k$.

\medskip\noindent
Assume $X$ is a proper curve over $k$.
By Divisors, Lemma \ref{divisors-lemma-filter-after-modification}
there exists a modification $f : X' \to X$ such that $f^*\mathcal{E}$
has a filtration whose successive quotients are invertible
$\mathcal{O}_{X'}$-modules. Since $f_*[X']_1 = [X]_1$ we conclude
from Lemma \ref{lemma-pushforward-cap-cj} that
$$
\deg(c_1(\mathcal{E}) \cap [X]_1) = \deg(c_1(f^*\mathcal{E}) \cap [X']_1)
$$
Since we have a similar relationship for the degree by
Varieties, Lemma \ref{varieties-lemma-degree-birational-pullback}
we reduce to the case where $\mathcal{E}$ has a filtration whose
successive quotients are invertible $\mathcal{O}_X$-modules.
In this case, we may use additivity of the degree
(Varieties, Lemma \ref{varieties-lemma-degree-additive})
and of first Chern classes (Lemma \ref{lemma-additivity-chern-classes})
to reduce to the case discussed in the next paragraph.

\medskip\noindent
Assume $X$ is a proper curve over $k$ and $\mathcal{E}$ is an
invertible $\mathcal{O}_X$-module. By
Divisors, Lemma
\ref{divisors-lemma-quasi-projective-Noetherian-pic-effective-Cartier}
we see that $\mathcal{E}$ is isomorphic to
$\mathcal{O}_X(D) \otimes \mathcal{O}_X(D')^{\otimes -1}$
for some effective Cartier divisors $D, D'$ on $X$ (this also uses
that $X$ is projective, see
Varieties, Lemma \ref{varieties-lemma-dim-1-proper-projective} for example).
By additivity of degree under tensor product of invertible sheaves
(Varieties, Lemma \ref{varieties-lemma-degree-tensor-product})
and additivity of $c_1$ under tensor product of invertible sheaves
(Lemma \ref{lemma-c1-cap-additive} or \ref{lemma-chern-classes-E-tensor-L})
we reduce to the case $\mathcal{E} = \mathcal{O}_X(D)$.
In this case the left hand side gives $\deg(D)$
(Varieties, Lemma \ref{varieties-lemma-degree-effective-Cartier-divisor})
and the right hand side gives $\deg([D]_0)$ by
Lemma \ref{lemma-geometric-cap}.
Since
$$
[D]_0 = \sum\nolimits_{x \in D}
\text{length}_{\mathcal{O}_{X, x}}(\mathcal{O}_{D, x}) [x] =
\sum\nolimits_{x \in D}
\text{length}_{\mathcal{O}_{D, x}}(\mathcal{O}_{D, x}) [x]
$$
by definition, we see
$$
\deg([D]_0) = \sum\nolimits_{x \in D}
\text{length}_{\mathcal{O}_{D, x}}(\mathcal{O}_{D, x}) [\kappa(x) : k] =
\dim_k \Gamma(D, \mathcal{O}_D) = \deg(D)
$$
The penultimate equality by
Algebra, Lemma \ref{algebra-lemma-pushdown-module}
using that $D$ is affine.
\end{proof}

\noindent
Finally, we can tie everything up with the numerical intersections
defined in Varieties, Section \ref{varieties-section-num}.

\begin{lemma}
\label{lemma-degrees-and-numerical-intersections}
Let $k$ be a field. Let $X$ be a proper scheme over $k$.
Let $Z \subset X$ be a closed subscheme of dimension $d$.
Let $\mathcal{L}_1, \ldots, \mathcal{L}_d$ be invertible
$\mathcal{O}_X$-modules. Then
$$
(\mathcal{L}_1 \cdots \mathcal{L}_d \cdot Z) =
\deg(
c_1(\mathcal{L}_1) \cap \ldots \cap c_1(\mathcal{L}_d) \cap [Z]_d)
$$
where the left hand side is defined in
Varieties, Definition \ref{varieties-definition-intersection-number}.
In particular,
$$
\deg_\mathcal{L}(Z) = \deg(c_1(\mathcal{L})^d \cap [Z]_d)
$$
if $\mathcal{L}$ is an ample invertible $\mathcal{O}_X$-module.
\end{lemma}

\begin{proof}
We will prove this by induction on $d$. If $d = 0$, then the result is
true by Varieties, Lemma \ref{varieties-lemma-chi-tensor-finite}.
Assume $d > 0$.

\medskip\noindent
Let $Z_i \subset Z$, $i = 1, \ldots, t$ be the irreducible components
of dimension $d$ with reduced induced scheme structure and let $m_i$ be the
multiplicity of $Z_i$ in $Z$. Then $[Z]_d = \sum m_i[Z_i]$ and
$c_1(\mathcal{L}_1) \cap \ldots \cap c_1(\mathcal{L}_d) \cap [Z]_d$
is the sum of the cycles
$m_i c_1(\mathcal{L}_1) \cap \ldots \cap c_1(\mathcal{L}_d) \cap [Z_i]$.
Since we have a similar decomposition for
$(\mathcal{L}_1 \cdots \mathcal{L}_d \cdot Z)$ by
Varieties, Lemma \ref{varieties-lemma-numerical-polynomial-leading-term}
it suffices to prove the lemma in case $Z = X$
is a proper variety of dimension $d$ over $k$.

\medskip\noindent
By Chow's lemma there exists a birational proper morphism $f : Y \to X$
with $Y$ H-projective over $k$. See Cohomology of Schemes, Lemma
\ref{coherent-lemma-chow-Noetherian} and Remark
\ref{coherent-remark-chow-Noetherian}. Then
$$
(f^*\mathcal{L}_1 \cdots f^*\mathcal{L}_d \cdot Y) =
(\mathcal{L}_1 \cdots \mathcal{L}_d \cdot X)
$$
by Varieties, Lemma \ref{varieties-lemma-intersection-number-and-pullback}
and we have
$$
f_*(c_1(f^*\mathcal{L}_1) \cap \ldots \cap c_1(f^*\mathcal{L}_d) \cap [Y]) =
c_1(\mathcal{L}_1) \cap \ldots \cap c_1(\mathcal{L}_d) \cap [X]
$$
by Lemma \ref{lemma-pushforward-cap-c1}. Thus we may replace $X$ by $Y$
and assume that $X$ is projective over $k$.

\medskip\noindent
If $X$ is a proper $d$-dimensional projective variety, then we can
write $\mathcal{L}_1 = \mathcal{O}_X(D) \otimes \mathcal{O}_X(D')^{\otimes -1}$
for some effective Cartier divisors $D, D' \subset X$
by Divisors, Lemma
\ref{divisors-lemma-quasi-projective-Noetherian-pic-effective-Cartier}.
By additivity for both sides of the equation
(Varieties, Lemma \ref{varieties-lemma-intersection-number-additive} and
Lemma \ref{lemma-c1-cap-additive})
we reduce to the case $\mathcal{L}_1 = \mathcal{O}_X(D)$ for some
effective Cartier divisor $D$.
By Varieties, Lemma
\ref{varieties-lemma-numerical-intersection-effective-Cartier-divisor}
we have
$$
(\mathcal{L}_1 \cdots \mathcal{L}_d \cdot X) =
(\mathcal{L}_2 \cdots \mathcal{L}_d \cdot D)
$$
and by Lemma \ref{lemma-geometric-cap} we have
$$
c_1(\mathcal{L}_1) \cap \ldots \cap c_1(\mathcal{L}_d) \cap [X] =
c_1(\mathcal{L}_2) \cap \ldots \cap c_1(\mathcal{L}_d) \cap [D]_{d - 1}
$$
Thus we obtain the result from our induction hypothesis.
\end{proof}








